% Creazione della copertina

\newcommand{\copertina}{
	\newgeometry{top=4cm}
	
	\begin{titlepage}
		\begin{center}
			
			\begin{center}
				%% qui metteteci l'immagine di copertina. Io ho messo quella dell'uni,
				%voi mettete quella del vostro grupo
				\centerline{\includegraphics[scale=0.8]{../logo}}
			\end{center}
			
			\vspace{1cm}
			
			\begin{Huge}
				\textbf{\Titolo{}} \\
			\end{Huge}
			
			\vspace{1cm}  
			
			\begin{large}
				\Gruppo{} \Data{}
			\end{large}	  
			
			\vspace{1cm}
			
			\bgroup
			\def\arraystretch{1.3}
			\centering
			\begin{tabular}{c|L{5cm}}
				\multicolumn{2}{c}{\textbf{About} } \\ \hline
			%	\textbf{Versione} &  \Versione{}\\
			%	\textbf{Responsabile} & \Responsabile{}\\
			%	\textbf{Verifica} & \Verifica{}\\
				\textbf{Company} & \Destinatoa{} \\
				\textbf{Author} & \Redazione{} \\
				\textbf{State} & \Stato{} \\
				\textbf{Use} & \Uso{} \\
				\textbf{Email} & \Email{} \\
			\end{tabular}
			\egroup
			
			\vspace{1.5cm}
			
			\begin{center}
				\textbf{Description\\}
				\DescrizioneDoc{}
			\end{center}
			
		\end{center}
	\end{titlepage}
	
	\restoregeometry
}

\definecolor{ao(english)}{rgb}{0.0, 0.5, 0.0}
\newcommand{\code}[1]{\textcolor{ao(english)}{\texttt{#1}}}

\newcommand{\gl}[1]{\textit{#1}\ped{g}}
